\paragraph{Explain working-capital movements from the verified history.}
The operating cycle now supports four reconciled movement schedules rather than
only its closing totals. \texttt{DeclaredWorkingCapital/1} first recomputes the
complete program result and selects one successful report, including the report
inside a provisional carry. Its actual opening book defines the period. Earlier
interim reports add no activity; a preceding carry establishes a new opening.
If execution later stops, an earlier successful report can still be selected,
while the failed command cannot stand in for a report.

For the illustrative cycle, receivables open at zero, increase by USD~50 of credit
sales and fall by USD~35 of collections, closing at USD~15. Inventory opens at
USD~20, receives USD~30 of purchases and releases USD~25 to cost of sales, closing
at USD~25. Supplier payables increase from zero by USD~30 and fall by USD~20;
wage payables increase by USD~10 and fall by USD~8. Their closing amounts are
USD~10 and USD~2. Every movement retains its event identity, date, typed inputs,
provenance, exact method definition and verified request/result references.
Liabilities use positive presentation even though the book stores credits as
negative balances. Purchases and deliveries also retain the supplied quantities
and the declared unit-cost premise.

Classification depends on the complete checked method definition, not its registry
nickname. Renaming a method key therefore preserves its meaning; changing the
method cannot masquerade as the old operation. An unsupported movement in a scoped
account prevents the schedule, even if another event offsets it. Gross posting
checks also detect an unsupported debit and credit that cancel within one event.
That last case exposed and repaired a first-draft error: inspecting only the net
journal had hidden the movement. The preserved counterexample distinguishes the
repair from a claim based only on successful ordinary examples.

After carry, collecting the remaining USD~15 and paying the remaining USD~10 of
supplier balances produces the next period's movements, with no repetition of
prior sales or purchases. Inventory remains USD~25 and wage payables USD~2.
The four working-capital changes explain a USD~5 cash contribution. In general,
this contribution is only part of an indirect cash reconciliation: a separate
noncash expense can leave a residual between profit, these four changes and direct
operating cash. The schedule displays that residual instead of forcing it to zero.

These schedules account for recorded movements. They retain unresolved receivable
ageing and impairment, inventory valuation and quantity coverage, maturity and
payment terms, related parties, tax and source-population completeness. An exact
roll-forward supplies useful note evidence, but cannot make those absent disclosures
complete. Both note-completeness and fiscal-completeness indicators remain false.

